\section{A successful extraction can still lose the fact}

Consider an Inline XBRL amount whose literal is \code{1,234}, whose admitted
numeric transform treats the comma as grouping, whose scale is $-3$, and whose
context and unit identify a particular dollar observation. The resulting amount
is $1234\times10^{-3}=1.234$. Scaling that normalized value again gives 0.001234
and is wrong. An adapter must retain the literal, transform, scale, context,
unit, nil status and original occurrence so that this error can be found.

The same principle applies to dates and dimensions. A lexical instant of
31 December remains distinct from the exclusive 1 January boundary used for
an interval. Publication and ingestion are separate timestamps. Two equal
numbers with different contexts remain separate occurrences. A nil fact is
unknown, not zero. Dropping a typed dimension or context reference yields a
preservation failure for that fact even if the resulting Markdown looks clear.

A PDF can preserve characters that the reader never sees. The retained HSBC
2023 annual report has a cover reading ``Annual Report and Accounts 2023'',
but both PyMuPDF and Poppler also return 2024 in the extracted cover text
\citep[cover, physical PDF p.~1]{hsbc2023annualcover}. Inspection of the drawing
order explains the discrepancy: the earlier 2024 text is covered by an opaque
dark rectangle, after which the 2023 text is painted. This is an observed
difference between extracted characters and the visible title. It does not
establish an accounting contradiction or explain why the file was constructed
that way. Agreement between two text extractors would not resolve it.

The proposed exchange keeps the original PDF, the unchanged raw extraction,
the page locator and a render from the same bytes. The annotation's
\code{literal} retains both years; separate \code{readings} record the raw
extraction and the proposed visible title. A linked inspection record identifies
the extractor versions and options, rendered image, relevant drawing order and
reviewer. Thus a raw-text search still finds 2024, while the inspected cover
supports a visible report-year reading of 2023. Rejecting 2024 as the visible
report year must not delete its extracted occurrence. The current example has
author inspection; independent adjudication remains pending.

\begingroup
\clubpenalty=10000
If an earlier review had already accepted 2024 as the report year, the correction
would require a new reading and an explicit \code{Supersede} decision. Results
depending on that reading would need revalidation in the new evidence view;
frozen earlier results would retain their original evidence. This follows the
review rules already stated above. Even after the visible title is resolved,
the cover year supplies no publication or availability timestamp for these
particular bytes. An as-of query must therefore retain the unresolved
availability condition. This single-page case motivates preservation and review
behavior; it supplies no extraction-accuracy ranking or general visibility
test for other pages.
\par\endgroup


The inspected Arelle interfaces support native XBRL contexts, dimensions,
transformations and validation. Its internal object IDs do not replace stable
source identity. The inspected Docling source offers useful layout and table
proposals, but its examined XBRL graph path does not retain every required
fact/context/unit/typed-dimension field. That path needs a preservation sidecar
or must report unsupported conversion for affected records. These are source
audit findings, not measured accuracy rankings of the tools.

A returned layout conversion can also conceal missing pages. At the pinned
Docling implementation, the timeout branch assigns \code{PARTIAL\_SUCCESS}
when pages were expected, even if none finished. The public \code{convert\_all}
interface permits that status through \code{raises\_on\_error=True} if the
remaining stages return normally. Thus absence of an exception cannot establish
that a three-page request produced any extracted content. The assembly code
can restore missing page numbers with empty page objects; counting those objects
does not count successful processing \citep{doclingpinned2026}. These are
source-derived failure cases, not observed failures on Apple's filing.

The proposed exchange therefore keeps the original page inventory, requested
range, successfully processed pages, empty placeholders and structured errors
separate. If only page two was requested and processed, that is success on the
declared subset; pages one and three remain outside the conversion claim. If
the three-page request times out before any page finishes, all three remain
unprocessed despite the partial-success status. Retain the source and the
unfinished pages for another attempt. A deadline is also distinct from stopped
execution: the inspected shutdown code waits for worker threads for a bounded
time and can report that a thread is still alive. Later execution must confirm
termination within its isolated process before treating the attempt as stopped
\citep{doclingpinned2026}.

A requested processing stage may also remain unperformed. In the same Docling
source, automatic OCR selects among engines according to platform, installed
packages, languages and runtime availability. If it selects no engine and has
recorded no language-specific rejection, it logs a warning; the subsequent call
returns the pages without OCR. Language-specific rejection can instead produce
an exception. Thus \code{do\_ocr=True} records a request, not evidence that text
recognition occurred. The proposed adapter records the requested mode, selected
engine and model identity, and processing outcome for each page. A page that
requires OCR but did not receive it remains unprocessed, even if other stages
produced layout regions \citep{doclingpinned2026}.

Identifying the engine still leaves its model to be fixed. The inspected default
layout preset names Heron at revision \code{main}, while the TableFormer
downloader names the model repository's \code{v2.3.0} tag. Neither name supplies
the hashes of the weights and preprocessing files actually used. The proposed
adapter binds those files, configurations and dependency versions before use.
Disabling remote inference services also does not disable model acquisition:
the layout helper downloads when no artifact directory is supplied. With a
supplied directory it requires the model's subfolder, but finding that folder
does not check its contents against the requested revision. Local file
verification and later isolated execution must therefore establish the intended
offline behavior; a source setting alone cannot do so \citep{doclingpinned2026}.

The text supplied to table recognition also needs an identity. The retained
\code{docling-ibm-models}~4.0.2 consumer uses OTSL cell and span tokens, converts
them to HTML, and merges text from separately supplied PDF cells
\citep{tableformerconsumer2026}. This is a later representation than the
original paper's HTML sequence. The Docling report explicitly describes a
subsequent token-language refinement \citep[Section 3]{nassar2024}; neither
that statement nor a model-card score identifies the later files with the
paper's experimental checkpoints.

The proposed adapter therefore retains the source page and crop, the supplied
text cells and their locations, the predicted table cells, and the assignments
and structural changes made during matching. If a correction combines columns
or adds a cell for unmatched text, a reviewer must be able to recover the
earlier structure and the text that prompted the change. Text obtained from an
OCR engine also retains that engine's identity and outcome. A resulting table
remains a proposal until its cell content, context and accounting interpretation
have been checked; a well-formed HTML table alone cannot settle those questions.


Token syntax and cell geometry require separate checks too. In the inspected
author version of OTSL, \citet[Sections 4.1--4.3]{lysak2023otsl} use five
tokens: C starts a cell, L joins the cell to its left, U joins the cell above,
X joins both, and NL ends a row. Neighbor restrictions and equal row lengths
allow invalid tokens to be detected during generation. The authors explicitly
distinguish a syntactically valid sequence from a correct prediction.

Even those local rules need a further condition to ensure rectangular merged
cells. Consider two rows, \code{C L NL} followed by \code{U C NL}. The L has
C to its left, the U has C above, and the first-row, first-column and
equal-length rules are satisfied; no X occurs. Yet the first cell joins the
top-left, top-right and bottom-left positions, leaving the bottom-right
position to the second C. It is L-shaped. Representing it with both a two-row
and a two-column HTML span would also occupy the second cell's position.
This finite example shows why the six printed rules alone do not establish a
nonoverlapping partition into rectangular cells. The proposed adapter checks
that each merged cell occupies a complete rectangle, that cells do not overlap,
and that every declared grid position belongs to a cell. Passing those checks
still leaves correspondence with the source table to be established.

The later decoder needs its own verification. In the retained 4.0.2 source,
a correction described as applying to the first row uses a row counter that
is initialized to zero and never updated. It therefore also rewrites an
affected token on later rows \citep{tableformerconsumer2026}. This is a
specific source-code defect relative to that stated scope; its frequency and
effect on actual model predictions have not been measured. Neither the paper's
syntax argument nor these partial corrections establish that a deployed
converter enforces all the proposed checks.


The retained Apple annual filing makes the dependency problem concrete. Its
issuer archive includes the instance, the issuer's schema and its calculation,
definition, label and presentation linkbases. Those files also refer to US GAAP
and SEC Reporting Taxonomy definitions, SEC schemas and XBRL definitions held
elsewhere. The issuer archive alone therefore cannot supply every definition
needed to interpret its tagged facts. For this particular archive, a static
inspection now resolves the recorded imports, link references and typed-domain
references to retained files, checks target namespaces and fragment identifiers,
and checks that instance facts have declared concepts, contexts and units.
This establishes a source dependency map. XML Schema and XBRL validation, the
issuer's interpretation of each concept, and the completeness of the wider
disclosure package remain separate questions.

An offline parser needs more than a map of intended filenames. At the inspected
Arelle source version, the cache lookup can return a local pathname when the
file is absent; subsequent loading records an \code{IOerror}. The proposed
adapter must check each mapped file and its hash before loading, retain the
missing dependency and affected records, and return an unresolved-dependency
result. It must not infer successful retrieval from a nonempty pathname or
silently substitute a newer taxonomy. Original schemas remain unchanged;
an external resolver or isolated cache associates their declared addresses
with the retained bytes. Any built-in cache fallback must belong to the same
checked source set \citep{arellepinned2026}.

Nor is a successful API return necessarily a validation verdict. In the same
source version, the controller's return value takes accumulated validation
errors into account only when \code{validationExitCode} is enabled. The proposed
adapter sets it explicitly, requests validation, disables automatic formula
execution, and examines the returned status, diagnostic messages and model
errors together. It keeps the model open while copying the necessary facts,
contexts, units and original occurrence references, then closes the session.
One session runs at a time in its owning thread within a process; the inspected
API uses shared package and plugin state. These are source-grounded interface
requirements for a later implementation, not results from an executed Arelle
adapter \citep{arellepinned2026}.

\input{sections/arelle-inline-values}

\input{sections/arelle-inline-values}

\input{sections/arelle-fact-identity}

Each adapter therefore pins code, dependency, model and taxonomy versions,
licenses, input and output schemas, transformations and field-loss behavior.
Preparation preserves bytes and receipts before parsing, inventories the
package, attaches occurrences and contexts, records normalization, checks
required fields independently and submits candidates for review. Partial
conversion retains both successful units and the exact failures or unexplored
frontier. Conversion counts cannot accept a package or an accounting treatment.

The proposed extraction comparison uses the same retained files, roots, budgets
and taxonomy resources. Manual located source units and transparent native
extraction provide the reference; applicable pdfplumber/Arelle configurations,
plain Docling and explicit context/footnote repairs form the candidate ladder.
Transcription, context, table structure, mapping, source support, retention,
historical admission and cost are separate outcomes. Lost hashes or locators,
look-ahead, silent omissions and untracked downloads veto the affected exchange.
The current 31 development annotations cannot establish a heldout accuracy
ranking, and cell counts within one report are not independent observations.

Reuse also has distribution conditions. Arelle's Apache and Docling's MIT code
licenses do not by themselves clear every model, taxonomy or dataset. ACTUS
technical-specification reuse retains its share-alike obligations; QuantLib,
Z3 and ledger tools retain their own code and dependency terms. The Beancount
and hledger comparison needs an explicit distribution boundary for copyleft
components. License admission records actual versions, notices, models and
data, rather than treating a repository's headline license as universal.

\enlargethispage{\baselineskip}
EasyOCR makes the dependency question concrete. Version~1.7.2 requires
\code{python-bidi} without a version constraint and imports its
\code{get\_display} function when the OCR module loads. Although the call
reorders Arabic recognition results, selecting English does not remove the
import or package requirement. In the inspected \code{python-bidi}~0.6.11
source, this import selects a Rust-backed implementation; the older Python
implementation has a different import path. Its own source states
LGPLv3-or-later terms. The inspected release is a candidate for dependency
analysis, and its source-level call signature alone does not establish
compatibility with an installed OCR configuration
\citep{easyocr172source,pythonbidi0611source}.

The same package shows why an accompanying notice still needs checking. Its
third-party list names older PyO3 and \code{unicode-bidi} versions than its
bundled Cargo lockfile. The exact locked sources retain MIT-or-Apache terms,
so the version discrepancy does not demonstrate a licence conflict. Other
locked packages add conditions: \code{target-lexicon} includes the LLVM
exception to Apache~2.0, while generated \code{unicode-ident} tables require
Unicode data-file terms alongside the chosen MIT or Apache licence
\citep{pythonbidi0611source,bidirustsources2026}. The adapter must therefore
bind the selected build and retain each applicable notice. Verified source
archives help prepare that decision; actual wheel contents, model rights
and the intended distribution still need separate assessment.

The actual taxonomy notices show why that distinction changes the design.
The inspected FASB notices permit specified financial reporting and unchanged
incorporation in works that explain or assist implementation, subject to their
notice and third-party conditions. They do not provide a general permission to
edit taxonomy definitions. Keeping the original files and recording location
mappings separately respects that distinction. The notices retrieved for this
inspection are unversioned; their relationship to the terms attached to the
2024 files still needs to be established before a distribution decision.
XBRL's public-document notice separately directs software and schemas to a
Software Notice, so permission for an explanatory document cannot simply be
applied to a schema bundle. Exact dependency licences, component-specific terms
and the intended distribution remain explicit preparation tasks
\citep{fasbtaxonomyuses,fasbsrtuses,xbrlcopyright}.
